%-------------------------------------------------------------------------------
%    WORK PLAN & PROGRESS PAGE (Шинэчлэгдсэн)
%-------------------------------------------------------------------------------

\begin{titlepage}

\vspace*{0.5cm}
Батлав. \deptname ын эрхлэгч: 
\begin{flushright}
\makebox[4cm]{\dotfill} /\chairname/ 
\end{flushright}

Удирдагч: 
\begin{flushright}
\makebox[4cm]{\dotfill} /\supname/
\end{flushright}

\begin{center}

\vspace*{1cm}
\textbf{{\large ТӨГСӨЛТИЙН АЖЛЫН ТӨЛӨВЛӨГӨӨ}}\\[0.5cm]
\textsc{\large Сэдэв: ''\ttitle''}\\[0.5cm]

\begin{tabular}{|c|p{7.5cm}|c|c|}
    \hline
    № & \makebox[7.5cm][c]{Ажлын бүлэг, хэсгийн нэр} & Эзлэх хувь & Дуусах хугацаа \\ \hline
    1 & \textbf{1-р үзлэг:} Удиртгал, Сэдвийн үндэслэл, Онолын судалгаа (CNN, RAG, MusicVAE) & 20\% & 2026-03-05 \\ \hline
    2 & \textbf{2-р үзлэг:} Системийн шинжилгээ, UML диаграммууд, Архитектур дизайн & 25\% & 9-р долоо хоног \\ \hline
    3 & \textbf{Төслийн хэсэг:} AI моделийн сургалт, Веб системийн хөгжүүлэлт, Интеграци & 40\% & 2026-05-01 \\ \hline
    4 & \textbf{Төгсгөл:} Туршилт, Үр дүнгийн шинжилгээ ба Дүгнэлт & 15\% & 2026-05-15 \\ \hline
\end{tabular}

\vspace{1.5cm}
Төлөвлөгөөг боловсруулсан оюутан: \makebox[3cm]{\dotfill} /\shortname/

\end{center}

\newpage

\begin{center}

\vspace*{1cm}
\textbf{{\large ТӨГСӨЛТИЙН АЖЛЫН ЯВЦ}}\\[0.5cm]

\begin{tabular}{|c|p{7.5cm}|c|c|}
    \hline
    № & \makebox[7.5cm][c]{Хийж гүйцэтгэсэн ажил} & Биелсэн хугацаа & Удирдагчийн гарын үсэг \\ \hline
    \textbf{1} & \textbf{1-р үзлэгийн бэлтгэл:} Сэдвийн онолын судалгаа, Удиртгал хэсэг бичигдсэн & \textbf{2026-03-04} & \\ \hline
    2 & 2-р үзлэгийн бэлтгэл: Системийн функциональ шаардлага, UML диаграммууд & 9-р долоо хоног & \\ \hline
    3 & AI Mixing Assistant болон RAG чатботын хөгжүүлэлт & - & \\ \hline
    4 & Marketplace болон Хэрэглэгчийн веб модуль & - & \\ \hline
\end{tabular}

\vspace{1cm}
\textbf{1-р үзлэгийн товч дүгнэлт:} \\[0.5cm]
\dotfill \\[0.2cm]
\dotfill \\[0.2cm]
\dotfill \\[0.5cm]
Удирдагч: \makebox[3cm]{\dotfill} /\supname/ \\

\vspace{1.5cm}
ЗӨВШӨӨРӨЛ \\[0.5cm]
Оюутан \shortname-ын бичсэн төгсөлтийн ажлыг УШК-д хамгаалуулахаар тодорхойлов.\\[0.5cm]
Салбарын эрхлэгч: \makebox[3cm]{\dotfill} /\chairname/
\end{center}

\end{titlepage}

\newpage

%-------------------------------------------------------------------------------
%    REVIEW PAGE (ШҮҮМЖ)
%-------------------------------------------------------------------------------

\begin{titlepage}
\begin{center}

{\scshape\Large \univname\par} 
{\scshape\large \facname\par}\vspace{1cm} 

\textbf{{\Large ШҮҮМЖИЙН ХУУДАС}}\\[1cm]

\end{center}

\normalsize

\deptname-ын салбарын төгсөх курсийн оюутан \shortname-ын ''\ttitle'' сэдэвт төгсөлтийн ажлын шүүмж.

\begin{enumerate}
\item Төслөөр дэвшүүлсэн асуудал (Хөгжим продюсингийн хүндрэл), AI болон DSP (Digital Signal Processing) онолын судалгаа:
\begin{center}
\dotfill \\[0.4cm]
\end{center}
\item Төслийн ерөнхий агуулга. CNN/MusicVAE загварын хэрэглээ, RAG чатботын бодит үр дүн:
\begin{center}
\dotfill \\[0.4cm]
\end{center}
\item Эмх цэгц, стандарт (LaTeX ашиглалт), зураг, хүснэгтийн тайлбар хангагдсан эсэх:
\begin{center}
\dotfill \\[0.4cm]
\end{center}
\item Төслөөр орхигдуулсан зүйлс болон цаашид анхаарах техникийн шаардлага:
\begin{center}
\dotfill \\[0.4cm]
\end{center}
\item Төслийн талаар онцолж тэмдэглэх зүйл (Инновацилаг байдал):
\begin{center}
\dotfill \\[0.4cm]
\end{center}
\item Ерөнхий оноо. (5 оноо)
\begin{center}
\dotfill \\[0.8cm]
\end{center}
\end{enumerate}

Шүүмж бичсэн: \makebox[3cm]{\dotfill} /\readname/ \\[0.5cm]
Ажлын газар: \dotfill \\[0.5cm]
Хаяг (Утас): \makebox[5cm]{\dotfill}
\end{titlepage}